\begin{figure}[H]
    \centering
    \definecolor{shellblue}{RGB}{80,145,195}
    \definecolor{pressurered}{RGB}{185,55,45}
    \begin{tikzpicture}[x=1cm,y=1cm,line cap=round,line join=round]
        \coordinate (O) at (3.1,2.4);
        \fill[shellblue!16] (O) circle (2.0);
        \fill[white] (O) circle (0.9);
        \draw[line width=1pt] (O) circle (2.0);
        \draw[line width=1pt] (O) circle (0.9);
        \fill (O) circle (0.035);
        \node[above left] at (O) {$z$};

        \draw[-latex,line width=0.8pt] (O) -- ++(35:0.9);
        \node[above right] at (3.65,2.95) {$a$};
        \draw[-latex,line width=0.8pt] (O) -- ++(-38:2.0);
        \node[below right] at (4.75,1.18) {$b$};

        \foreach \angle in {35,115,195,275} {
            \draw[pressurered,-latex,line width=0.9pt]
                ($(O)+(\angle:0.88)$) -- ($(O)+(\angle:1.30)$);
        }
        \node[pressurered,align=left] at (6.1,3.35) {presión interna\\uniforme $p$};
        \draw[pressurered,line width=0.6pt] (5.25,3.2) -- (4.1,2.95);
        \node[align=left] at (6.2,1.35) {pared externa\\sin tracción};
        \draw[black!55,line width=0.6pt] (5.3,1.55) -- (4.85,1.15);
        \node[align=center] at (8.0,2.4) {sección transversal\\del cilindro hueco};
    \end{tikzpicture}
    \caption{Geometría del cilindro de paredes gruesas y presión en su superficie interna.}
\end{figure}